\begin{abstract}
The square-safe Sieve Sequence certifies that every survivor
$p\in[Q,Q^2)$ is prime. Requiring $p+2$ to survive the same filters would ask
for a twin-prime pair. We study a deliberately weaker target instead:
require $p+2$ to avoid primes below $z=Q^{2\alpha}$ for some fixed
$\alpha>1/3$. Positivity then implies that $p+2$ has at most two prime
factors for all sufficiently large $Q$.

We prove three exact algebraic results for this relaxed weight. The
one-divisor count has an explicit local factor and one periodic boundary
remainder. The natural pre-sieved divisor remainder is exactly a discrepancy
of primes in the progression $-2$ modulo $d$. Finally, the scalar-centered
final weight decomposes into nonprincipal character products
$\chi(m)\chi(n)$. A quadratic character modulo $3$ correlates with the full
relaxed survivor count on the complete reduced wheel, refuting the shortcut
that scalar-density centering creates arbitrary-coefficient Type-II
orthogonality.

These results do not prove positivity of the relaxed weight. They identify its
correct Type-I comparison, refute one over-strong Type-II formulation, and
isolate the remaining need for an averaged prime-progression theorem followed
by a locally adapted bilinear estimate.
\end{abstract}
